\documentclass[10pt,a4paper,twocolumn]{article}

% =========================================================
% Language and typography
% =========================================================
\usepackage[utf8]{inputenc}
\usepackage[T1]{fontenc}
\usepackage[english,french]{babel}
\usepackage{lmodern}
\usepackage{microtype}

% =========================================================
% Page layout - COMPACT
% =========================================================
\usepackage[a4paper,margin=0.8cm,top=1cm,bottom=0.8cm,columnsep=0.6cm]{geometry}
\usepackage{setspace}
\setstretch{0.85}

% Remove paragraph indent and spacing
\usepackage{parskip}
% \setlength{\parindent}{0pt}
% \setlength{\parskip}{0.15em}

% Reduce spacing around sections
\usepackage{titlesec}
% \titlespacing*{\section}{0pt}{0.2em}{0.1em}
% \titlespacing*{\subsection}{0pt}{0.15em}{0.08em}
% \titlespacing*{\subsubsection}{0pt}{0.1em}{0.05em}
% \titleformat{\section}{\normalfont\bfseries\footnotesize}{\thesection}{0.4em}{}
% \titleformat{\subsection}{\normalfont\bfseries\small}{\thesubsection}{0.3em}{}
% \titleformat{\subsubsection}{\normalfont\itshape\small}{\thesubsubsection}{0.2em}{}

% Reduce spacing in lists
\usepackage{enumitem}
\setlist{nolistsep, leftmargin=0.6em, itemsep=0.02em}

% =========================================================
% Math and symbols
% =========================================================
\usepackage{amsmath,amssymb,amsfonts}

% =========================================================
% Tables and figures
% =========================================================
\usepackage{graphicx}
\usepackage{booktabs}
\usepackage{tabularx}
\usepackage{array}
\usepackage{xcolor}
\usepackage{tikz}
\usepackage{float}

% =========================================================
% Links
% =========================================================
\usepackage[hidelinks]{hyperref}

% =========================================================
% Boxes - COMPACT
% =========================================================
\usepackage[most]{tcolorbox}

\definecolor{SpaceBlue}{RGB}{18,35,60}
\definecolor{SpaceCyan}{RGB}{0,150,190}
\definecolor{LightBlue}{RGB}{235,245,250}
\definecolor{LightGray}{RGB}{245,245,245}
\definecolor{WarningOrange}{RGB}{230,130,40}

\tcbset{
    colback=LightGray,
    colframe=SpaceBlue,
    arc=0.8mm,
    boxrule=0.2pt,
    left=0.4mm,
    right=0.4mm,
    top=0.2mm,
    bottom=0.2mm,
    fonttitle=\bfseries\footnotesize
}

\newtcolorbox{keybox}{
    title=Key idea,
    colback=LightBlue,
    colframe=SpaceCyan
}

\newtcolorbox{warningbox}{
    title=Pitfall,
    colback=orange!7,
    colframe=WarningOrange
}

\newtcolorbox{vocabbox}{
    title=Vocabulary,
    colback=white,
    colframe=SpaceBlue
}

\newtcolorbox{questionbox}{
    title=Q\&A,
    colback=white,
    colframe=black!60
}

% =========================================================
% Custom commands
% =========================================================

\newcommand{\keyword}[1]{\textbf{\textcolor{SpaceBlue}{#1}}}
\newcommand{\eng}[1]{\textit{#1}}
\newcommand{\fr}[1]{\textit{#1}}

\newcommand{\lesson}[3]{
\section{#1}
\noindent
\textbf{Goal:} #2 \newline
\textbf{Why:} #3
\vspace{0.1em}
}

% =========================================================
% Title
% =========================================================

\title{\tiny \textbf{MASSpace 2026 Cheatsheet} \\ Space Systems, Autonomy, and Multi-Agent Systems}
\author{\tiny Julien Soulé}
\date{\tiny \today}

% =========================================================
% Document
% =========================================================

\begin{document}

\selectlanguage{english}

\maketitle

% {\tiny \tableofcontents}

% =========================================================
% Preface
% =========================================================


% \section*{How to Use These Notes}
% \addcontentsline{toc}{section}{How to Use These Notes}

% This document is designed as a compact revision guide for MASSpace 2026.
% Each lesson is written as a concise revision sheet, with:
% \begin{itemize}
%     \item core concepts;
%     \item English technical vocabulary;
%     \item operational intuition;
%     \item links with Multi-Agent Systems and MARL;
%     \item typical questions and answers.
% \end{itemize}

% The goal is not to become an aerospace engineer, but to become fluent enough to discuss space autonomy, satellite operations, and multi-agent coordination with specialists.

% \newpage

% =========================================================
% PART 1
% =========================================================

\part{Foundations for Non-Space Researchers}

% =========================================================
% INTRODUCTION
% =========================================================

\section*{Part Overview}

This first part introduces the minimum space-systems culture required to comfortably discuss:
\begin{itemize}
    \item satellite missions;
    \item operational constraints;
    \item orbital dynamics;
    \item communication limitations;
    \item distributed autonomy in space systems.
\end{itemize}

The objective is not to become an orbital mechanics expert, but to develop:
\begin{itemize}
    \item operational intuition;
    \item technical vocabulary;
    \item systems-level understanding;
    \item the ability to follow scientific discussions at MASSpace.
\end{itemize}

% \newpage

% =========================================================
% LESSON 1
% =========================================================

\lesson
{The Space Domain: Big Picture}
{Understand the major categories of space missions, actors, and operational ecosystems.}
{Space autonomy research is deeply connected to mission operations, satellite constellations, and distributed decision-making.}

% ---------------------------------------------------------
% WHY SPACE MATTERS
% ---------------------------------------------------------

\subsection{Why Space Systems Matter}

Modern societies increasingly rely on space infrastructures:
\begin{itemize}
    \item GPS and navigation;
    \item Earth observation;
    \item weather forecasting;
    \item telecommunications;
    \item military intelligence;
    \item disaster monitoring;
    \item scientific exploration.
\end{itemize}

Space systems are becoming:
\begin{itemize}
    \item larger;
    \item more distributed;
    \item more autonomous;
    \item more software-defined;
    \item more AI-driven.
\end{itemize}

This creates major challenges in:
\begin{itemize}
    \item coordination;
    \item resilience;
    \item scheduling;
    \item autonomy;
    \item cybersecurity;
    \item mission supervision.
\end{itemize}

\begin{keybox}
    Future space systems are increasingly viewed as distributed cyber-physical multi-agent systems.
\end{keybox}

% ---------------------------------------------------------
% MAIN TYPES OF MISSIONS
% ---------------------------------------------------------

\subsection{Main Types of Space Missions}

\subsubsection*{1. Earth Observation (EO)}

Goal:
observe Earth using optical, radar, or hyperspectral sensors.

Typical applications:
\begin{itemize}
    \item climate monitoring;
    \item agriculture;
    \item defense intelligence;
    \item disaster response;
    \item maritime surveillance.
\end{itemize}

Examples:
\begin{itemize}
    \item Sentinel constellation;
    \item Planet Labs;
    \item Maxar;
    \item ICEYE.
\end{itemize}

\begin{vocabbox}
    \begin{tabularx}{\columnwidth}{lX}
        \toprule
        English term           & Meaning                                      \\
        \midrule
        Earth Observation (EO) & Observation of Earth from space              \\
        Imaging satellite      & Satellite carrying optical/radar sensors     \\
        Remote sensing         & Acquiring information remotely using sensors \\
        Revisit time           & Time between two observations of same area   \\
        Ground resolution      & Precision of captured imagery                \\
        \bottomrule
    \end{tabularx}
\end{vocabbox}

% ---------------------------------------------------------

\subsubsection*{2. Telecommunications}

Goal:
relay information globally.

Applications:
\begin{itemize}
    \item internet access;
    \item TV broadcasting;
    \item secure military communications;
    \item IoT communications.
\end{itemize}

Examples:
\begin{itemize}
    \item Starlink;
    \item OneWeb;
    \item Iridium.
\end{itemize}

\begin{keybox}
    Mega-constellations massively increase coordination and scheduling complexity.
\end{keybox}

% ---------------------------------------------------------

\subsubsection*{3. Navigation Systems}

Goal:
provide precise positioning and timing services.

Examples:
\begin{itemize}
    \item GPS (USA);
    \item Galileo (EU);
    \item GLONASS (Russia);
    \item BeiDou (China).
\end{itemize}

Applications:
\begin{itemize}
    \item transportation;
    \item military operations;
    \item banking synchronization;
    \item autonomous vehicles.
\end{itemize}

% ---------------------------------------------------------

\subsubsection*{4. Scientific Missions}

Goal:
study:
\begin{itemize}
    \item planets;
    \item stars;
    \item black holes;
    \item cosmology;
    \item astrobiology.
\end{itemize}

Examples:
\begin{itemize}
    \item James Webb Space Telescope;
    \item Hubble;
    \item Voyager;
    \item Mars rovers.
\end{itemize}

% ---------------------------------------------------------

\subsubsection*{5. Defense and Dual-Use Systems}

Increasingly important topic.

Applications:
\begin{itemize}
    \item reconnaissance;
    \item missile warning;
    \item secure communications;
    \item electronic warfare;
    \item space situational awareness.
\end{itemize}

\begin{warningbox}
    Many modern space technologies are dual-use:
    they have both civilian and military applications.
\end{warningbox}

% ---------------------------------------------------------
% SPACE ECOSYSTEM
% ---------------------------------------------------------

\subsection{The Space Ecosystem}

Space systems involve several interacting layers.

\begin{center}
    Launch segment
    $\rightarrow$
    Space segment
    $\rightarrow$
    Ground segment
    $\rightarrow$
    Users
\end{center}

% ---------------------------------------------------------

\subsubsection*{1. Space Segment}

Includes:
\begin{itemize}
    \item satellites;
    \item spacecraft;
    \item orbital infrastructures.
\end{itemize}

This is where onboard autonomy executes.

% ---------------------------------------------------------

\subsubsection*{2. Ground Segment}

Includes:
\begin{itemize}
    \item ground stations;
    \item mission control centers;
    \item telemetry processing systems;
    \item scheduling systems.
\end{itemize}

Ground operators:
\begin{itemize}
    \item monitor satellites;
    \item send commands;
    \item analyze mission data;
    \item manage anomalies.
\end{itemize}

% ---------------------------------------------------------

\subsubsection*{3. User Segment}

Final users:
\begin{itemize}
    \item governments;
    \item civilians;
    \item companies;
    \item defense organizations.
\end{itemize}

% ---------------------------------------------------------
% CONSTELLATIONS
% ---------------------------------------------------------

\subsection{Constellations}

A constellation is a coordinated group of satellites.

Examples:
\begin{itemize}
    \item Starlink;
    \item OneWeb;
    \item Planet Labs.
\end{itemize}

Constellations are interesting for MAS research because they involve:
\begin{itemize}
    \item decentralized agents;
    \item local observations;
    \item communication constraints;
    \item distributed scheduling;
    \item cooperative objectives;
    \item dynamic environments.
\end{itemize}

\begin{keybox}
    Satellite constellations are naturally modeled as distributed multi-agent systems.
\end{keybox}

% ---------------------------------------------------------
% LINK WITH MARL
% ---------------------------------------------------------

\subsection{Connection with MARL and MAS}

Most modern space autonomy problems can be formulated as:
\begin{itemize}
    \item Dec-POMDPs;
    \item distributed optimization problems;
    \item multi-agent coordination problems;
    \item constrained planning problems.
\end{itemize}

Typical MARL challenges:
\begin{itemize}
    \item partial observability;
    \item sparse communication;
    \item dynamic task allocation;
    \item delayed rewards;
    \item safety constraints;
    \item non-stationarity.
\end{itemize}

% ---------------------------------------------------------
% TYPICAL DISCUSSION
% ---------------------------------------------------------

\subsection{Typical Scientific Discussion}

\begin{questionbox}
    \textbf{Q: Why is autonomy becoming necessary in space systems?}

    A: Because future constellations are becoming too large, dynamic, and communication-constrained for purely manual centralized operations.
\end{questionbox}

\begin{questionbox}
    \textbf{Q: Why are constellations relevant for MAS research?}

    A: Because they involve distributed agents coordinating under partial observability, communication constraints, and shared mission objectives.
\end{questionbox}

% \newpage

% =========================================================
% LESSON 2
% =========================================================

\lesson
{Basic Orbital Mechanics for AI Researchers}
{Build operational intuition about orbital motion and visibility constraints.}
{Orbital mechanics directly influence communication opportunities, scheduling, and coordination.}

% ---------------------------------------------------------
% BIG PICTURE
% ---------------------------------------------------------

\subsection{Why Orbital Mechanics Matter}

Orbital mechanics determine:
\begin{itemize}
    \item where satellites are;
    \item when they can observe targets;
    \item when they can communicate;
    \item when conjunction risks appear;
    \item how often regions can be revisited.
\end{itemize}

For AI researchers:
\begin{itemize}
    \item you do NOT need advanced astrodynamics;
    \item you DO need operational intuition.
\end{itemize}

% ---------------------------------------------------------
% BASIC INTUITION
% ---------------------------------------------------------

\subsection{Basic Intuition}

A satellite stays in orbit because:
\begin{itemize}
    \item gravity pulls it toward Earth;
    \item its velocity continuously makes it ``fall around'' Earth.
\end{itemize}

\begin{keybox}
    Orbit = continuous free fall around Earth.
\end{keybox}

Higher orbit:
\begin{itemize}
    \item slower orbital speed;
    \item larger coverage;
    \item longer orbital period.
\end{itemize}

Lower orbit:
\begin{itemize}
    \item faster orbit;
    \item better resolution;
    \item lower communication latency.
\end{itemize}

% ---------------------------------------------------------
% MAIN ORBITS
% ---------------------------------------------------------

\subsection{Main Orbit Families}

% ---------------------------------------------------------

\subsubsection*{1. Low Earth Orbit (LEO)}

Altitude:
\begin{itemize}
    \item roughly 160--2000 km.
\end{itemize}

Characteristics:
\begin{itemize}
    \item low latency;
    \item high revisit dynamics;
    \item small coverage area;
    \item many satellites required.
\end{itemize}

Used for:
\begin{itemize}
    \item Earth observation;
    \item mega-constellations;
    \item low-latency communications.
\end{itemize}

\begin{vocabbox}
    \begin{tabularx}{\columnwidth}{lX}
        \toprule
        English term & Meaning                               \\
        \midrule
        LEO          & Low Earth Orbit                       \\
        Revisit time & Time before observing same area again \\
        Orbital pass & Satellite passing over region/station \\
        Low latency  & Small communication delay             \\
        \bottomrule
    \end{tabularx}
\end{vocabbox}

% ---------------------------------------------------------

\subsubsection*{2. Medium Earth Orbit (MEO)}

Altitude:
\begin{itemize}
    \item 2000--35000 km.
\end{itemize}

Mostly used for:
\begin{itemize}
    \item navigation systems.
\end{itemize}

Examples:
\begin{itemize}
    \item GPS;
    \item Galileo.
\end{itemize}

% ---------------------------------------------------------

\subsubsection*{3. Geostationary Orbit (GEO)}

Altitude:
\begin{itemize}
    \item approximately 35786 km.
\end{itemize}

Characteristics:
\begin{itemize}
    \item satellite appears fixed relative to Earth;
    \item huge coverage area;
    \item high latency.
\end{itemize}

Used for:
\begin{itemize}
    \item telecommunications;
    \item weather satellites.
\end{itemize}

% ---------------------------------------------------------
% IMPORTANT ORBITAL CONCEPTS
% ---------------------------------------------------------

\subsection{Important Orbital Concepts}

% ---------------------------------------------------------

\subsubsection*{Inclination}

Inclination defines orbit tilt relative to Earth equator.

Examples:
\begin{itemize}
    \item equatorial orbit;
    \item polar orbit;
    \item sun-synchronous orbit.
\end{itemize}

Polar orbits are important for:
\begin{itemize}
    \item global Earth observation.
\end{itemize}

% ---------------------------------------------------------

\subsubsection*{Orbital Period}

Time needed to complete one orbit.

Typical LEO:
\begin{itemize}
    \item about 90 minutes.
\end{itemize}

% ---------------------------------------------------------

\subsubsection*{Ground Track}

Trajectory projected on Earth surface.

Determines:
\begin{itemize}
    \item revisit opportunities;
    \item coverage dynamics.
\end{itemize}

% ---------------------------------------------------------
% VISIBILITY
% ---------------------------------------------------------

\subsection{Visibility and Communication}

Communication depends on:
\begin{itemize}
    \item orbital geometry;
    \item line of sight;
    \item antenna orientation;
    \item Earth obstruction.
\end{itemize}

\begin{keybox}
    Space communication is naturally intermittent.
\end{keybox}

This is extremely important for:
\begin{itemize}
    \item decentralized autonomy;
    \item onboard planning;
    \item distributed coordination;
    \item resilient operations.
\end{itemize}

% ---------------------------------------------------------
% MARL CONNECTION
% ---------------------------------------------------------

\subsection{Connection with MARL}

Orbital dynamics create:
\begin{itemize}
    \item time-varying connectivity graphs;
    \item delayed coordination opportunities;
    \item dynamic observation windows;
    \item non-stationary interaction structures.
\end{itemize}

This naturally connects with:
\begin{itemize}
    \item Dec-POMDPs;
    \item graph-based MARL;
    \item communication learning;
    \item scheduling under uncertainty.
\end{itemize}

% ---------------------------------------------------------
% TYPICAL DISCUSSION
% ---------------------------------------------------------

\subsection{Typical Scientific Discussion}

\begin{questionbox}
    \textbf{Q: Why are communication disruptions common in LEO constellations?}

    A: Because satellites move quickly relative to Earth and to each other, constantly changing visibility conditions and communication opportunities.
\end{questionbox}

\begin{questionbox}
    \textbf{Q: Why is decentralized autonomy important in LEO systems?}

    A: Because communication opportunities are intermittent and centralized control may not scale operationally.
\end{questionbox}

% \newpage

% =========================================================
% LESSON 3
% =========================================================

\lesson
{The Space Environment}
{Understand environmental hazards and operational degradation sources in space.}
{Future autonomous systems must remain robust under harsh and uncertain conditions.}

% ---------------------------------------------------------
% INTRO
% ---------------------------------------------------------

\subsection{Why the Space Environment is Difficult}

Space is:
\begin{itemize}
    \item hostile;
    \item resource-constrained;
    \item highly dynamic;
    \item difficult to repair physically.
\end{itemize}

Unlike terrestrial systems:
\begin{itemize}
    \item hardware replacement is difficult;
    \item communication delays may be large;
    \item failures can become mission-ending.
\end{itemize}

\begin{warningbox}
    Space systems are safety-critical and mission-critical systems.
\end{warningbox}

% ---------------------------------------------------------
% SPACE DEBRIS
% ---------------------------------------------------------

\subsection{Space Debris}

Space debris includes:
\begin{itemize}
    \item inactive satellites;
    \item rocket fragments;
    \item collision fragments;
    \item abandoned objects.
\end{itemize}

Debris creates:
\begin{itemize}
    \item conjunction risks;
    \item collision threats;
    \item operational complexity.
\end{itemize}

% ---------------------------------------------------------

\subsubsection*{Conjunction Assessment}

A conjunction:
\begin{itemize}
    \item is a close approach between orbital objects.
\end{itemize}

Operators must:
\begin{itemize}
    \item estimate collision probability;
    \item decide whether avoidance maneuvers are necessary.
\end{itemize}

\begin{vocabbox}
    \begin{tabularx}{\columnwidth}{lX}
        \toprule
        English term                 & Meaning                                  \\
        \midrule
        Space debris                 & Orbital junk and inactive objects        \\
        Conjunction                  & Close orbital encounter                  \\
        Collision avoidance maneuver & Orbit adjustment reducing collision risk \\
        SSA / SDA                    & Space Situational / Domain Awareness     \\
        \bottomrule
    \end{tabularx}
\end{vocabbox}

% ---------------------------------------------------------
% RADIATION
% ---------------------------------------------------------

\subsection{Radiation}

Space radiation can:
\begin{itemize}
    \item damage electronics;
    \item corrupt memory;
    \item create transient faults.
\end{itemize}

Important concept:
\begin{itemize}
    \item Single-Event Upset (SEU).
\end{itemize}

An SEU:
\begin{itemize}
    \item is a bit flip caused by radiation.
\end{itemize}

This matters for:
\begin{itemize}
    \item onboard AI;
    \item autonomous decision-making;
    \item fault tolerance.
\end{itemize}

% ---------------------------------------------------------
% SOLAR WEATHER
% ---------------------------------------------------------

\subsection{Solar Weather}

The Sun can generate:
\begin{itemize}
    \item solar flares;
    \item geomagnetic storms;
    \item energetic particle events.
\end{itemize}

Effects:
\begin{itemize}
    \item communication disruption;
    \item navigation degradation;
    \item electronic failures.
\end{itemize}

% ---------------------------------------------------------
% COMMUNICATION CONSTRAINTS
% ---------------------------------------------------------

\subsection{Communication Constraints}

Space systems often experience:
\begin{itemize}
    \item intermittent connectivity;
    \item latency;
    \item bandwidth limitations;
    \item temporary isolation.
\end{itemize}

This strongly motivates:
\begin{itemize}
    \item onboard autonomy;
    \item decentralized execution;
    \item resilient coordination.
\end{itemize}

% ---------------------------------------------------------
% CYBER AND ADVERSARIAL
% ---------------------------------------------------------

\subsection{Cyber and Adversarial Threats}

Modern space systems may operate in:
\begin{itemize}
    \item contested environments;
    \item adversarial conditions;
    \item cyber-degraded situations.
\end{itemize}

Typical threats:
\begin{itemize}
    \item jamming;
    \item spoofing;
    \item denial of service;
    \item malware insertion.
\end{itemize}

\begin{keybox}
    Future resilient autonomy must handle both physical and cyber degradation.
\end{keybox}

% ---------------------------------------------------------
% MARL CONNECTION
% ---------------------------------------------------------

\subsection{Connection with MARL and MAS}

Environmental hazards naturally create:
\begin{itemize}
    \item uncertainty;
    \item stochastic transitions;
    \item degraded observations;
    \item constrained coordination.
\end{itemize}

Interesting research directions:
\begin{itemize}
    \item robust MARL;
    \item safe RL;
    \item resilient coordination;
    \item fault-aware planning;
    \item adaptive autonomy.
\end{itemize}

% ---------------------------------------------------------
% TYPICAL DISCUSSION
% ---------------------------------------------------------

\subsection{Typical Scientific Discussion}

\begin{questionbox}
    \textbf{Q: Why is resilience especially important in space systems?}

    A: Because space systems operate in harsh environments where failures are difficult or impossible to repair physically.
\end{questionbox}

\begin{questionbox}
    \textbf{Q: Why is decentralized autonomy important for resilience?}

    A: Because communication may be disrupted or delayed, requiring satellites to continue operating locally under degraded conditions.
\end{questionbox}


% =========================================================
% PART 2
% =========================================================

\part{Space Systems Engineering}

% =========================================================
% INTRODUCTION
% =========================================================

\section*{Part Overview}

This part introduces the engineering foundations of space systems.

The objective is to understand:
\begin{itemize}
    \item how satellites are designed;
    \item how they communicate;
    \item how missions are operated;
    \item which operational constraints drive autonomy research.
\end{itemize}

For AI and MAS researchers, this part is extremely important because:
\begin{itemize}
    \item autonomy is constrained by hardware and physics;
    \item mission operations involve scheduling and coordination;
    \item communication is intermittent and resource-constrained;
    \item satellites are cyber-physical systems.
\end{itemize}

\begin{keybox}
    Space autonomy problems are fundamentally systems-engineering problems.
\end{keybox}

% \newpage

% =========================================================
% LESSON 1
% =========================================================

\lesson
{Satellite Architecture}
{Understand the major subsystems composing a satellite.}
{Autonomous decision-making depends on energy, communication, orientation, onboard computing, and mission payload constraints.}

% ---------------------------------------------------------
% BIG PICTURE
% ---------------------------------------------------------

\subsection{What is a Satellite?}

A satellite is:
\begin{itemize}
    \item a spacecraft orbiting a celestial body;
    \item composed of several interacting subsystems;
    \item designed to execute one or several missions.
\end{itemize}

A satellite is essentially:
\begin{center}
    payload + platform
\end{center}

Where:
\begin{itemize}
    \item the payload performs the mission;
    \item the platform supports the payload operationally.
\end{itemize}

Examples:
\begin{itemize}
    \item imaging payload;
    \item radar payload;
    \item communication payload;
    \item scientific instruments.
\end{itemize}

% ---------------------------------------------------------
% MAIN SUBSYSTEMS
% ---------------------------------------------------------

\subsection{Main Satellite Subsystems}

% ---------------------------------------------------------

\subsubsection*{1. Payload}

The payload performs the actual mission.

Examples:
\begin{itemize}
    \item optical camera;
    \item radar;
    \item hyperspectral sensor;
    \item communication antenna.
\end{itemize}

Payload operations are often:
\begin{itemize}
    \item energy-intensive;
    \item storage-intensive;
    \item scheduling-constrained.
\end{itemize}

\begin{keybox}
    In many autonomy problems, payload usage is the primary mission objective.
\end{keybox}

% ---------------------------------------------------------

\subsubsection*{2. Power Subsystem}

Responsible for:
\begin{itemize}
    \item energy generation;
    \item battery storage;
    \item power distribution.
\end{itemize}

Typical components:
\begin{itemize}
    \item solar panels;
    \item batteries;
    \item power regulators.
\end{itemize}

Energy is one of the most important constraints in autonomous operations.

Examples:
\begin{itemize}
    \item imaging consumes energy;
    \item communication consumes energy;
    \item attitude maneuvers consume energy.
\end{itemize}

% ---------------------------------------------------------

\subsubsection*{3. Attitude Determination and Control System (ADCS)}

ADCS controls:
\begin{itemize}
    \item satellite orientation;
    \item stabilization;
    \item pointing accuracy.
\end{itemize}

Important because:
\begin{itemize}
    \item sensors must point toward targets;
    \item antennas must point toward communication partners;
    \item solar panels may require optimal orientation.
\end{itemize}

Typical actuators:
\begin{itemize}
    \item reaction wheels;
    \item thrusters;
    \item magnetorquers.
\end{itemize}

\begin{vocabbox}
    \begin{tabularx}{\columnwidth}{lX}
        \toprule
        English term  & Meaning                                    \\
        \midrule
        Attitude      & Orientation of spacecraft                  \\
        Pointing      & Directing instrument/antenna toward target \\
        Slew maneuver & Rotational reorientation maneuver          \\
        Stabilization & Maintaining desired orientation            \\
        \bottomrule
    \end{tabularx}
\end{vocabbox}

% ---------------------------------------------------------

\subsubsection*{4. Communication Subsystem}

Responsible for:
\begin{itemize}
    \item uplink;
    \item downlink;
    \item inter-satellite communications.
\end{itemize}

Communication constraints strongly influence:
\begin{itemize}
    \item distributed autonomy;
    \item coordination;
    \item scheduling;
    \item fault tolerance.
\end{itemize}

% ---------------------------------------------------------

\subsubsection*{5. Onboard Computer (OBC)}

The onboard computer:
\begin{itemize}
    \item executes software;
    \item manages autonomy;
    \item processes telemetry;
    \item controls payloads and subsystems.
\end{itemize}

Important limitation:
\begin{itemize}
    \item space hardware is often computationally constrained.
\end{itemize}

Why?
\begin{itemize}
    \item radiation-hardened processors are less powerful;
    \item reliability is prioritized over raw performance.
\end{itemize}

\begin{warningbox}
    Space AI systems often run on hardware much weaker than modern GPUs.
\end{warningbox}

% ---------------------------------------------------------

\subsubsection*{6. Thermal Control}

Spacecraft experience:
\begin{itemize}
    \item extreme heating;
    \item extreme cooling.
\end{itemize}

Thermal systems regulate temperature using:
\begin{itemize}
    \item radiators;
    \item insulation;
    \item heaters.
\end{itemize}

Thermal constraints may influence:
\begin{itemize}
    \item payload scheduling;
    \item operational planning;
    \item safe autonomy.
\end{itemize}

% ---------------------------------------------------------

\subsubsection*{7. Propulsion}

Used for:
\begin{itemize}
    \item orbit corrections;
    \item collision avoidance;
    \item station keeping;
    \item formation flying.
\end{itemize}

Propellant is limited:
\begin{itemize}
    \item maneuvers must be optimized carefully.
\end{itemize}

% ---------------------------------------------------------
% SYSTEMS THINKING
% ---------------------------------------------------------

\subsection{Systems Thinking}

Satellite operations involve strong subsystem couplings.

Example:
\begin{center}
    Observation task
    $\rightarrow$
    attitude maneuver
    $\rightarrow$
    energy consumption
    $\rightarrow$
    communication delay
    $\rightarrow$
    mission scheduling impact
\end{center}

\begin{keybox}
    Space systems are highly coupled cyber-physical systems.
\end{keybox}

% ---------------------------------------------------------
% CONNECTION WITH AI
% ---------------------------------------------------------

\subsection{Connection with AI and MARL}

Autonomous agents must reason under:
\begin{itemize}
    \item energy constraints;
    \item communication limitations;
    \item orientation constraints;
    \item limited onboard computing;
    \item safety requirements.
\end{itemize}

Interesting research topics:
\begin{itemize}
    \item resource-aware planning;
    \item constrained RL;
    \item mission scheduling;
    \item safe autonomy;
    \item adaptive mission management.
\end{itemize}

% ---------------------------------------------------------
% TYPICAL DISCUSSION
% ---------------------------------------------------------

\subsection{Typical Scientific Discussion}

\begin{questionbox}
    \textbf{Q: Why are onboard AI systems constrained in space?}

    A: Because space-qualified hardware prioritizes reliability and radiation resistance over computational performance.
\end{questionbox}

\begin{questionbox}
    \textbf{Q: Why is energy management critical in autonomous satellite operations?}

    A: Because most actions consume energy, and satellites operate with limited power budgets and battery capacity.
\end{questionbox}

% \newpage

% =========================================================
% LESSON 2
% =========================================================

\lesson
{Space Communications}
{Understand how communication constraints shape space autonomy.}
{Communication limitations are one of the primary motivations for decentralized autonomy.}

% ---------------------------------------------------------
% BIG PICTURE
% ---------------------------------------------------------

\subsection{Why Space Communications are Difficult}

Unlike terrestrial networks:
\begin{itemize}
    \item satellites move constantly;
    \item line of sight changes over time;
    \item bandwidth is limited;
    \item latency may be significant;
    \item communication windows are intermittent.
\end{itemize}

Consequences:
\begin{itemize}
    \item centralized control becomes difficult;
    \item local autonomy becomes necessary;
    \item coordination opportunities become sparse.
\end{itemize}

% ---------------------------------------------------------
% MAIN COMMUNICATION TYPES
% ---------------------------------------------------------

\subsection{Main Communication Types}

% ---------------------------------------------------------

\subsubsection*{1. Uplink}

Ground
$\rightarrow$
Satellite.

Used for:
\begin{itemize}
    \item telecommands;
    \item software updates;
    \item mission reconfiguration.
\end{itemize}

% ---------------------------------------------------------

\subsubsection*{2. Downlink}

Satellite
$\rightarrow$
Ground.

Used for:
\begin{itemize}
    \item telemetry;
    \item scientific data;
    \item imagery;
    \item mission logs.
\end{itemize}

% ---------------------------------------------------------

\subsubsection*{3. Inter-Satellite Links (ISL)}

Satellite
$\leftrightarrow$
Satellite.

Used for:
\begin{itemize}
    \item distributed coordination;
    \item relaying;
    \item collaborative autonomy.
\end{itemize}

Very important for:
\begin{itemize}
    \item mega-constellations;
    \item swarm autonomy;
    \item distributed planning.
\end{itemize}

% ---------------------------------------------------------
% COMMUNICATION WINDOWS
% ---------------------------------------------------------

\subsection{Communication Windows}

Communication depends on:
\begin{itemize}
    \item orbital geometry;
    \item visibility;
    \item antenna orientation;
    \item Earth obstruction.
\end{itemize}

This creates:
\begin{itemize}
    \item intermittent communication;
    \item delayed synchronization;
    \item dynamic network topologies.
\end{itemize}

\begin{keybox}
    Connectivity in space systems is often time-varying and opportunistic.
\end{keybox}

% ---------------------------------------------------------
% LATENCY
% ---------------------------------------------------------

\subsection{Latency}

Latency:
\begin{itemize}
    \item is the delay before information is received.
\end{itemize}

Higher orbit:
\begin{itemize}
    \item larger latency.
\end{itemize}

Important consequences:
\begin{itemize}
    \item slower centralized control;
    \item delayed replanning;
    \item asynchronous coordination.
\end{itemize}

% ---------------------------------------------------------
% DTN
% ---------------------------------------------------------

\subsection{Delay-Tolerant Networking (DTN)}

DTN:
\begin{itemize}
    \item is designed for intermittent communication environments.
\end{itemize}

Core idea:
\begin{itemize}
    \item store data locally;
    \item forward opportunistically later.
\end{itemize}

This resembles:
\begin{itemize}
    \item distributed buffering;
    \item asynchronous relaying;
    \item opportunistic routing.
\end{itemize}

Very relevant for:
\begin{itemize}
    \item ORBITAL-like environments;
    \item relay missions;
    \item decentralized scheduling.
\end{itemize}

% ---------------------------------------------------------
% GROUND STATIONS
% ---------------------------------------------------------

\subsection{Ground Stations}

Ground stations:
\begin{itemize}
    \item are limited resources;
    \item must be scheduled carefully.
\end{itemize}

Challenges:
\begin{itemize}
    \item many satellites;
    \item few stations;
    \item dynamic priorities.
\end{itemize}

This creates:
\begin{itemize}
    \item scheduling problems;
    \item resource allocation problems;
    \item coordination challenges.
\end{itemize}

% ---------------------------------------------------------
% CONNECTION WITH MAS
% ---------------------------------------------------------

\subsection{Connection with MAS and MARL}

Communication constraints naturally create:
\begin{itemize}
    \item partial observability;
    \item asynchronous coordination;
    \item delayed information sharing;
    \item decentralized planning requirements.
\end{itemize}

Interesting research directions:
\begin{itemize}
    \item communication-aware MARL;
    \item graph neural coordination;
    \item adaptive communication policies;
    \item distributed scheduling.
\end{itemize}

% ---------------------------------------------------------
% TYPICAL DISCUSSION
% ---------------------------------------------------------

\subsection{Typical Scientific Discussion}

\begin{questionbox}
    \textbf{Q: Why is centralized control difficult in large constellations?}

    A: Because communication opportunities are intermittent, latency can be high, and the number of satellites becomes difficult to manage operationally.
\end{questionbox}

\begin{questionbox}
    \textbf{Q: Why are inter-satellite links important?}

    A: Because they enable distributed coordination and decentralized data relaying without relying entirely on ground infrastructure.
\end{questionbox}

% \newpage

% =========================================================
% LESSON 3
% =========================================================

\lesson
{Mission Operations}
{Understand how space missions are operated and supervised.}
{Mission operations are the real-world layer where autonomy, planning, and coordination become operationally useful.}

% ---------------------------------------------------------
% BIG PICTURE
% ---------------------------------------------------------

\subsection{What are Mission Operations?}

Mission operations involve:
\begin{itemize}
    \item monitoring spacecraft;
    \item planning activities;
    \item scheduling observations;
    \item managing anomalies;
    \item coordinating resources.
\end{itemize}

Mission operations centers:
\begin{itemize}
    \item supervise satellite fleets;
    \item send commands;
    \item analyze telemetry;
    \item ensure mission continuity.
\end{itemize}

% ---------------------------------------------------------
% TELEMETRY
% ---------------------------------------------------------

\subsection{Telemetry and Telecommands}

% ---------------------------------------------------------

\subsubsection*{Telemetry}

Telemetry:
\begin{itemize}
    \item is data sent from spacecraft to ground.
\end{itemize}

Examples:
\begin{itemize}
    \item battery level;
    \item temperature;
    \item payload status;
    \item position;
    \item fault reports.
\end{itemize}

% ---------------------------------------------------------

\subsubsection*{Telecommands}

Telecommands:
\begin{itemize}
    \item are instructions sent from ground to spacecraft.
\end{itemize}

Examples:
\begin{itemize}
    \item activate payload;
    \item rotate spacecraft;
    \item upload mission plan;
    \item execute maneuver.
\end{itemize}

% ---------------------------------------------------------
% MISSION PLANNING
% ---------------------------------------------------------

\subsection{Mission Planning}

Mission planning determines:
\begin{itemize}
    \item what should be done;
    \item when;
    \item by which satellite;
    \item under which constraints.
\end{itemize}

Planning constraints:
\begin{itemize}
    \item energy;
    \item visibility;
    \item communication windows;
    \item pointing constraints;
    \item thermal constraints;
    \item mission priorities.
\end{itemize}

\begin{keybox}
    Mission planning is fundamentally a constrained optimization and scheduling problem.
\end{keybox}

% ---------------------------------------------------------
% SCHEDULING
% ---------------------------------------------------------

\subsection{Scheduling}

Scheduling:
\begin{itemize}
    \item allocates resources over time.
\end{itemize}

Examples:
\begin{itemize}
    \item observation scheduling;
    \item downlink scheduling;
    \item maneuver scheduling;
    \item ground station scheduling.
\end{itemize}

Important properties:
\begin{itemize}
    \item combinatorial complexity;
    \item uncertainty;
    \item dynamic priorities;
    \item real-time adaptation.
\end{itemize}

% ---------------------------------------------------------
% ANOMALIES
% ---------------------------------------------------------

\subsection{Anomaly and Fault Management}

Satellites experience:
\begin{itemize}
    \item hardware faults;
    \item communication losses;
    \item software anomalies;
    \item environmental perturbations.
\end{itemize}

Mission operations must:
\begin{itemize}
    \item detect anomalies;
    \item diagnose failures;
    \item recover mission capability.
\end{itemize}

This is often called:
\begin{itemize}
    \item FDIR:
          Fault Detection, Isolation, and Recovery.
\end{itemize}

% ---------------------------------------------------------
% AUTONOMY LEVELS
% ---------------------------------------------------------

\subsection{Autonomy Levels}

Space missions can operate at different autonomy levels.

% ---------------------------------------------------------

\subsubsection*{Low Autonomy}

Ground operators:
\begin{itemize}
    \item make most decisions manually.
\end{itemize}

% ---------------------------------------------------------

\subsubsection*{Medium Autonomy}

Satellite:
\begin{itemize}
    \item executes predefined procedures autonomously.
\end{itemize}

Ground:
\begin{itemize}
    \item supervises globally.
\end{itemize}

% ---------------------------------------------------------

\subsubsection*{High Autonomy}

Satellite:
\begin{itemize}
    \item performs local planning;
    \item adapts dynamically;
    \item coordinates with peers.
\end{itemize}

Very important for:
\begin{itemize}
    \item mega-constellations;
    \item deep-space missions;
    \item degraded environments.
\end{itemize}

% ---------------------------------------------------------
% HUMAN-IN-THE-LOOP
% ---------------------------------------------------------

\subsection{Human-in-the-Loop Operations}

Modern autonomy research usually keeps:
\begin{itemize}
    \item human supervision;
    \item mission-level control;
    \item operator validation.
\end{itemize}

Why?
\begin{itemize}
    \item safety;
    \item trust;
    \item explainability;
    \item certification;
    \item accountability.
\end{itemize}

\begin{keybox}
    Most realistic operational visions involve supervised autonomy rather than full autonomous replacement.
\end{keybox}

% ---------------------------------------------------------
% CONNECTION WITH AI
% ---------------------------------------------------------

\subsection{Connection with AI and MAS}

Mission operations strongly connect with:
\begin{itemize}
    \item planning;
    \item scheduling;
    \item multi-agent coordination;
    \item constrained optimization;
    \item reinforcement learning;
    \item explainable autonomy.
\end{itemize}

Interesting research directions:
\begin{itemize}
    \item autonomous mission orchestration;
    \item adaptive scheduling;
    \item resilient planning;
    \item organizational autonomy;
    \item mission-aware MARL.
\end{itemize}

% ---------------------------------------------------------
% TYPICAL DISCUSSION
% ---------------------------------------------------------

\subsection{Typical Scientific Discussion}

\begin{questionbox}
    \textbf{Q: Why is scheduling difficult in space missions?}

    A: Because many constraints interact simultaneously: visibility, energy, communication opportunities, thermal limits, priorities, and uncertainty.
\end{questionbox}

\begin{questionbox}
    \textbf{Q: Why is human supervision still important in autonomous missions?}

    A: Because space systems are mission-critical and safety-critical, requiring trust, explainability, accountability, and operational validation.
\end{questionbox}

% =========================================================
% PART 3
% =========================================================

\part{Space Autonomy and AI}

% =========================================================
% INTRODUCTION
% =========================================================

\section*{Part Overview}

This part focuses on autonomous decision-making in space systems.

The objective is to understand:
\begin{itemize}
    \item why autonomy is becoming necessary;
    \item how AI can support mission operations;
    \item why decentralized coordination matters;
    \item how MAS and MARL connect naturally with space systems.
\end{itemize}

This is the core conceptual bridge between:
\begin{itemize}
    \item aerospace systems;
    \item distributed systems;
    \item multi-agent systems;
    \item reinforcement learning;
    \item resilient autonomy.
\end{itemize}

\begin{keybox}
    Future space systems are increasingly evolving from remotely controlled infrastructures toward autonomous distributed cyber-physical systems.
\end{keybox}

% \newpage

% =========================================================
% LESSON 1
% =========================================================

\lesson
{Autonomous Space Systems}
{Understand why autonomy is becoming central in modern space missions.}
{Future constellations are too large, dynamic, and communication-constrained for purely manual centralized operations.}

% ---------------------------------------------------------
% WHY AUTONOMY
% ---------------------------------------------------------

\subsection{Why Space Autonomy Matters}

Traditional missions relied heavily on:
\begin{itemize}
    \item centralized control;
    \item predefined procedures;
    \item human operators.
\end{itemize}

However, future missions involve:
\begin{itemize}
    \item larger constellations;
    \item dynamic operational contexts;
    \item communication disruptions;
    \item onboard decision-making;
    \item adversarial environments.
\end{itemize}

This creates strong pressure toward:
\begin{itemize}
    \item autonomous operations;
    \item decentralized coordination;
    \item adaptive planning.
\end{itemize}

\begin{keybox}
    Autonomy is becoming necessary primarily because of operational scalability limitations.
\end{keybox}

% ---------------------------------------------------------
% MAIN AUTONOMY FUNCTIONS
% ---------------------------------------------------------

\subsection{Main Functions of Space Autonomy}

% ---------------------------------------------------------

\subsubsection*{1. Autonomous Planning}

The spacecraft:
\begin{itemize}
    \item generates plans locally;
    \item adapts to new events;
    \item prioritizes actions dynamically.
\end{itemize}

Examples:
\begin{itemize}
    \item selecting observation targets;
    \item choosing relay opportunities;
    \item managing energy usage.
\end{itemize}

% ---------------------------------------------------------

\subsubsection*{2. Replanning}

Unexpected events may occur:
\begin{itemize}
    \item communication loss;
    \item sensor failure;
    \item urgent observation request;
    \item collision alert.
\end{itemize}

The system must:
\begin{itemize}
    \item revise plans dynamically;
    \item preserve mission continuity.
\end{itemize}

% ---------------------------------------------------------

\subsubsection*{3. Fault Management}

Autonomous systems must:
\begin{itemize}
    \item detect anomalies;
    \item isolate failures;
    \item recover safely.
\end{itemize}

This is often called:
\begin{center}
    FDIR
\end{center}

Fault Detection, Isolation, and Recovery.

% ---------------------------------------------------------

\subsubsection*{4. Resource Management}

Autonomy must manage:
\begin{itemize}
    \item energy;
    \item communication bandwidth;
    \item onboard storage;
    \item thermal constraints;
    \item maneuver budgets.
\end{itemize}

\begin{warningbox}
    Most autonomy decisions in space are constrained-resource decisions.
\end{warningbox}

% ---------------------------------------------------------
% LEVELS OF AUTONOMY
% ---------------------------------------------------------

\subsection{Levels of Autonomy}

% ---------------------------------------------------------

\subsubsection*{Low Autonomy}

Human operators:
\begin{itemize}
    \item make most operational decisions.
\end{itemize}

Satellite:
\begin{itemize}
    \item executes predefined procedures.
\end{itemize}

% ---------------------------------------------------------

\subsubsection*{Medium Autonomy}

Satellite:
\begin{itemize}
    \item executes local reactions autonomously;
    \item handles simple contingencies.
\end{itemize}

Ground:
\begin{itemize}
    \item supervises mission globally.
\end{itemize}

% ---------------------------------------------------------

\subsubsection*{High Autonomy}

Satellite:
\begin{itemize}
    \item plans dynamically;
    \item coordinates with peers;
    \item adapts to uncertainty;
    \item manages degraded conditions.
\end{itemize}

Important for:
\begin{itemize}
    \item mega-constellations;
    \item swarm systems;
    \item deep-space missions.
\end{itemize}

% ---------------------------------------------------------
% HUMAN SUPERVISION
% ---------------------------------------------------------

\subsection{Human Supervision}

Modern operational visions usually keep:
\begin{itemize}
    \item operators in the loop;
    \item mission-level supervision;
    \item human validation authority.
\end{itemize}

This is called:
\begin{center}
    Human-in-the-loop autonomy
\end{center}

or:
\begin{center}
    Supervisory autonomy
\end{center}

\begin{keybox}
    Most realistic space autonomy architectures augment operators rather than replace them.
\end{keybox}

% ---------------------------------------------------------
% CONNECTION WITH AI
% ---------------------------------------------------------

\subsection{Connection with AI and MARL}

Autonomous operations naturally involve:
\begin{itemize}
    \item sequential decision-making;
    \item uncertainty;
    \item partial observability;
    \item dynamic environments;
    \item distributed coordination.
\end{itemize}

This strongly connects with:
\begin{itemize}
    \item reinforcement learning;
    \item planning;
    \item Dec-POMDPs;
    \item constrained optimization;
    \item organizational MAS.
\end{itemize}

% ---------------------------------------------------------
% TYPICAL DISCUSSION
% ---------------------------------------------------------

\subsection{Typical Scientific Discussion}

\begin{questionbox}
    \textbf{Q: Why is autonomy becoming unavoidable in future constellations?}

    A: Because constellation scale and operational complexity are increasing faster than human operators can manage manually.
\end{questionbox}

\begin{questionbox}
    \textbf{Q: Why is full autonomy still difficult operationally?}

    A: Because space systems require trust, explainability, safety guarantees, certification, and robust behavior under degraded conditions.
\end{questionbox}

% \newpage

% =========================================================
% LESSON 2
% =========================================================

\lesson
{Multi-Agent Systems in Space}
{Understand why space systems naturally map to MAS problems.}
{Satellite constellations are inherently distributed systems requiring coordination under uncertainty.}

% ---------------------------------------------------------
% WHY MAS
% ---------------------------------------------------------

\subsection{Why Space Systems are Multi-Agent Systems}

Modern constellations involve:
\begin{itemize}
    \item multiple satellites;
    \item distributed observations;
    \item decentralized resources;
    \item dynamic communication opportunities.
\end{itemize}

Each satellite:
\begin{itemize}
    \item has local observations;
    \item has local constraints;
    \item performs local decisions;
    \item contributes to global mission objectives.
\end{itemize}

\begin{keybox}
    A satellite constellation is naturally a distributed multi-agent cyber-physical system.
\end{keybox}

% ---------------------------------------------------------
% MAIN MAS PROBLEMS
% ---------------------------------------------------------

\subsection{Main MAS Problems in Space}

% ---------------------------------------------------------

\subsubsection*{1. Distributed Task Allocation}

Question:
\begin{center}
    Which satellite should perform which task?
\end{center}

Examples:
\begin{itemize}
    \item Earth observation requests;
    \item relaying tasks;
    \item collision monitoring;
    \item communication routing.
\end{itemize}

Challenges:
\begin{itemize}
    \item limited resources;
    \item visibility constraints;
    \item dynamic priorities.
\end{itemize}

% ---------------------------------------------------------

\subsubsection*{2. Distributed Scheduling}

Question:
\begin{center}
    When should tasks be executed?
\end{center}

Must consider:
\begin{itemize}
    \item communication windows;
    \item energy levels;
    \item payload availability;
    \item thermal constraints.
\end{itemize}

% ---------------------------------------------------------

\subsubsection*{3. Coordination}

Satellites may need to:
\begin{itemize}
    \item cooperate;
    \item relay data;
    \item share information;
    \item synchronize actions.
\end{itemize}

Examples:
\begin{itemize}
    \item collaborative imaging;
    \item relay chains;
    \item distributed sensing.
\end{itemize}

% ---------------------------------------------------------

\subsubsection*{4. Consensus and Agreement}

Agents may need to agree on:
\begin{itemize}
    \item priorities;
    \item task ownership;
    \item mission states;
    \item distributed plans.
\end{itemize}

% ---------------------------------------------------------
% SWARM SYSTEMS
% ---------------------------------------------------------

\subsection{Swarm and Formation Systems}

Swarm systems involve:
\begin{itemize}
    \item large numbers of cooperative agents;
    \item decentralized behaviors;
    \item emergent coordination.
\end{itemize}

Formation flying:
\begin{itemize}
    \item satellites maintain geometric configurations.
\end{itemize}

Applications:
\begin{itemize}
    \item distributed sensing;
    \item interferometry;
    \item Earth observation;
    \item adaptive communication networks.
\end{itemize}

% ---------------------------------------------------------
% DEC-POMDP
% ---------------------------------------------------------

\subsection{Connection with Dec-POMDPs}

Space systems naturally satisfy Dec-POMDP characteristics:
\begin{itemize}
    \item decentralized agents;
    \item partial observability;
    \item stochastic environment;
    \item delayed rewards;
    \item shared objectives.
\end{itemize}

Typical observations:
\begin{itemize}
    \item local energy state;
    \item local visibility;
    \item communication availability;
    \item nearby agents.
\end{itemize}

Typical actions:
\begin{itemize}
    \item observe;
    \item relay;
    \item maneuver;
    \item transmit;
    \item idle.
\end{itemize}

% ---------------------------------------------------------
% ORGANIZATIONAL MAS
% ---------------------------------------------------------

\subsection{Organizational MAS in Space}

Organizational approaches introduce:
\begin{itemize}
    \item explicit roles;
    \item missions;
    \item constraints;
    \item coordination semantics.
\end{itemize}

Examples:
\begin{itemize}
    \item observer satellites;
    \item relay satellites;
    \item safety-oriented agents.
\end{itemize}

Advantages:
\begin{itemize}
    \item improved explainability;
    \item controllability;
    \item operational discipline.
\end{itemize}

% ---------------------------------------------------------
% TYPICAL DISCUSSION
% ---------------------------------------------------------

\subsection{Typical Scientific Discussion}

\begin{questionbox}
    \textbf{Q: Why are Dec-POMDPs relevant for constellations?}

    A: Because satellites operate with local observations, limited communication, uncertainty, and shared mission objectives.
\end{questionbox}

\begin{questionbox}
    \textbf{Q: Why are organizational approaches interesting for space autonomy?}

    A: Because they provide explicit coordination structure and improve controllability and explainability in distributed autonomous systems.
\end{questionbox}

% \newpage

% =========================================================
% LESSON 3
% =========================================================

\lesson
{AI for Space Operations}
{Understand how AI can support satellite operations and mission management.}
{AI can improve adaptability and scalability, but operational deployment introduces strong safety and explainability requirements.}

% ---------------------------------------------------------
% WHY AI
% ---------------------------------------------------------

\subsection{Why AI is Becoming Important}

Space operations are becoming:
\begin{itemize}
    \item more dynamic;
    \item more distributed;
    \item more data-intensive;
    \item more resource-constrained.
\end{itemize}

AI can help:
\begin{itemize}
    \item automate scheduling;
    \item optimize coordination;
    \item detect anomalies;
    \item support operators;
    \item improve resilience.
\end{itemize}

\begin{keybox}
    The main operational motivation for AI is scalability under uncertainty.
\end{keybox}

% ---------------------------------------------------------
% ONBOARD AI
% ---------------------------------------------------------

\subsection{Onboard AI}

Onboard AI:
\begin{itemize}
    \item executes directly on spacecraft hardware.
\end{itemize}

Advantages:
\begin{itemize}
    \item faster reactions;
    \item reduced communication dependence;
    \item local adaptation.
\end{itemize}

Challenges:
\begin{itemize}
    \item limited hardware;
    \item energy constraints;
    \item safety requirements;
    \item radiation robustness.
\end{itemize}

% ---------------------------------------------------------
% EDGE AI
% ---------------------------------------------------------

\subsection{Edge AI in Space}

Edge AI:
\begin{itemize}
    \item processes data near acquisition sources;
    \item reduces downlink requirements;
    \item supports distributed autonomy.
\end{itemize}

Examples:
\begin{itemize}
    \item onboard image filtering;
    \item autonomous target detection;
    \item adaptive communication routing.
\end{itemize}

% ---------------------------------------------------------
% RL APPLICATIONS
% ---------------------------------------------------------

\subsection{Reinforcement Learning Applications}

RL can support:
\begin{itemize}
    \item task scheduling;
    \item observation prioritization;
    \item communication routing;
    \item adaptive planning;
    \item resource management.
\end{itemize}

Interesting RL properties:
\begin{itemize}
    \item adaptation;
    \item sequential reasoning;
    \item long-term optimization.
\end{itemize}

% ---------------------------------------------------------
% LIMITATIONS OF PURE RL
% ---------------------------------------------------------

\subsection{Limitations of Pure RL}

Operational deployment remains difficult because RL policies may:
\begin{itemize}
    \item lack explainability;
    \item violate operational doctrine;
    \item behave unpredictably;
    \item fail under distribution shifts.
\end{itemize}

This motivates:
\begin{itemize}
    \item safe RL;
    \item constrained RL;
    \item organizational RL;
    \item neuro-symbolic autonomy.
\end{itemize}

\begin{warningbox}
    High reward does not necessarily imply operational acceptability.
\end{warningbox}

% ---------------------------------------------------------
% EXPLAINABILITY
% ---------------------------------------------------------

\subsection{Explainability and Auditability}

Operators need to understand:
\begin{itemize}
    \item why decisions are made;
    \item whether policies remain safe;
    \item how coordination emerges.
\end{itemize}

Important concepts:
\begin{itemize}
    \item explainable autonomy;
    \item interpretable coordination;
    \item trajectory-level analysis;
    \item operational auditability.
\end{itemize}

% ---------------------------------------------------------
% HUMAN-AI TEAMING
% ---------------------------------------------------------

\subsection{Human-Autonomy Teaming}

Modern visions usually involve:
\begin{itemize}
    \item AI-assisted operations;
    \item operator supervision;
    \item mixed-initiative planning;
    \item adjustable autonomy.
\end{itemize}

Important concept:
\begin{center}
    Operational assist
\end{center}

Meaning:
\begin{itemize}
    \item autonomy supports operators rather than replacing them entirely.
\end{itemize}

% ---------------------------------------------------------
% CONNECTION WITH YOUR RESEARCH
% ---------------------------------------------------------

\subsection{Connection with Organizational MARL}

Interesting research directions:
\begin{itemize}
    \item organizational guidance;
    \item role-based coordination;
    \item constrained autonomy;
    \item mission-aware policies;
    \item explainable MARL;
    \item neuro-symbolic world models.
\end{itemize}

These directions attempt to combine:
\begin{itemize}
    \item adaptation;
    \item controllability;
    \item robustness;
    \item operational interpretability.
\end{itemize}

% ---------------------------------------------------------
% TYPICAL DISCUSSION
% ---------------------------------------------------------

\subsection{Typical Scientific Discussion}

\begin{questionbox}
    \textbf{Q: Why is pure RL difficult to deploy operationally in space systems?}

    A: Because operators require safety guarantees, explainability, predictability, and robustness under degraded conditions.
\end{questionbox}

\begin{questionbox}
    \textbf{Q: Why is onboard AI useful?}

    A: Because communication opportunities may be limited, delayed, or disrupted, requiring local decision-making capabilities.
\end{questionbox}

% \newpage

% =========================================================
% LESSON 4
% =========================================================

\lesson
{Space Cybersecurity and Resilience}
{Understand cyber threats and resilience challenges in space systems.}
{Future autonomous constellations may operate in adversarial and contested environments.}

% ---------------------------------------------------------
% WHY CYBERSECURITY MATTERS
% ---------------------------------------------------------

\subsection{Why Cybersecurity Matters in Space}

Modern satellites are:
\begin{itemize}
    \item software-defined;
    \item network-connected;
    \item increasingly autonomous.
\end{itemize}

Consequences:
\begin{itemize}
    \item cyberattacks become possible;
    \item disruptions may propagate across constellations;
    \item autonomy itself becomes an attack surface.
\end{itemize}

% ---------------------------------------------------------
% MAIN THREATS
% ---------------------------------------------------------

\subsection{Main Threats}

% ---------------------------------------------------------

\subsubsection*{1. Jamming}

Jamming:
\begin{itemize}
    \item intentionally disrupts communication signals.
\end{itemize}

Effects:
\begin{itemize}
    \item communication loss;
    \item degraded coordination;
    \item reduced situational awareness.
\end{itemize}

% ---------------------------------------------------------

\subsubsection*{2. Spoofing}

Spoofing:
\begin{itemize}
    \item injects fake signals or information.
\end{itemize}

Examples:
\begin{itemize}
    \item fake GPS signals;
    \item false telemetry;
    \item deceptive commands.
\end{itemize}

% ---------------------------------------------------------

\subsubsection*{3. Malware and Intrusions}

Potential targets:
\begin{itemize}
    \item onboard software;
    \item communication infrastructure;
    \item ground stations.
\end{itemize}

% ---------------------------------------------------------
% RESILIENCE
% ---------------------------------------------------------

\subsection{Resilience}

Resilience:
\begin{itemize}
    \item is the ability to preserve mission capability under failures or attacks.
\end{itemize}

Important concepts:
\begin{itemize}
    \item graceful degradation;
    \item fault tolerance;
    \item adaptive autonomy;
    \item distributed recovery.
\end{itemize}

\begin{keybox}
    A resilient system does not necessarily preserve full performance. It preserves mission value under degradation.
\end{keybox}

% ---------------------------------------------------------
% CONTESTED ENVIRONMENTS
% ---------------------------------------------------------

\subsection{Contested Environments}

Future systems may operate in:
\begin{itemize}
    \item adversarial environments;
    \item degraded communication contexts;
    \item cyber-contested domains.
\end{itemize}

This strongly motivates:
\begin{itemize}
    \item decentralized autonomy;
    \item local decision-making;
    \item robust coordination.
\end{itemize}

% ---------------------------------------------------------
% CONNECTION WITH MARL
% ---------------------------------------------------------

\subsection{Connection with MARL and MAS}

Interesting research directions:
\begin{itemize}
    \item robust MARL;
    \item adversarial MARL;
    \item cyber-resilient coordination;
    \item distributed recovery strategies;
    \item graceful degradation policies.
\end{itemize}

% ---------------------------------------------------------
% TYPICAL DISCUSSION
% ---------------------------------------------------------

\subsection{Typical Scientific Discussion}

\begin{questionbox}
    \textbf{Q: Why is resilience important in future constellations?}

    A: Because large distributed systems are likely to experience failures, communication disruptions, and potentially adversarial conditions.
\end{questionbox}

\begin{questionbox}
    \textbf{Q: Why does distributed autonomy help resilience?}

    A: Because local agents can continue operating even when centralized coordination or communication is degraded.
\end{questionbox}



% =========================================================
% PART 4
% =========================================================

\part{Advanced Research Topics}

% =========================================================
% INTRODUCTION
% =========================================================

\section*{Part Overview}

This part focuses on frontier research topics at the intersection of:
\begin{itemize}
    \item space systems;
    \item artificial intelligence;
    \item autonomy;
    \item resilience;
    \item multi-agent coordination.
\end{itemize}

The objective is not only to understand the current state of the field, but also:
\begin{itemize}
    \item identify open research challenges;
    \item understand emerging scientific directions;
    \item connect your own research with future MASSpace topics.
\end{itemize}

This part is particularly important for:
\begin{itemize}
    \item workshop discussions;
    \item networking;
    \item identifying collaboration opportunities;
    \item proposing original research ideas.
\end{itemize}

\begin{keybox}
    Modern space autonomy research increasingly focuses on scalable, resilient, explainable, and distributed autonomous systems.
\end{keybox}

% \newpage

% =========================================================
% LESSON 1
% =========================================================

\lesson
{Resilient Multi-Agent Space Systems}
{Understand how distributed autonomous systems can remain operational under degradation and failures.}
{Future constellations must preserve mission capability despite uncertainty, faults, and adversarial conditions.}

% ---------------------------------------------------------
% WHY RESILIENCE
% ---------------------------------------------------------

\subsection{Why Resilience Matters}

Space systems operate in:
\begin{itemize}
    \item harsh environments;
    \item uncertain conditions;
    \item communication-constrained settings;
    \item potentially adversarial contexts.
\end{itemize}

Failures are unavoidable:
\begin{itemize}
    \item hardware degradation;
    \item communication disruptions;
    \item software faults;
    \item cyberattacks;
    \item unexpected environmental events.
\end{itemize}

The question is therefore not:
\begin{center}
    ``Can failures occur?''
\end{center}

But:
\begin{center}
    ``How can mission value be preserved under degradation?''
\end{center}

% ---------------------------------------------------------
% GRACEFUL DEGRADATION
% ---------------------------------------------------------

\subsection{Graceful Degradation}

Graceful degradation:
\begin{itemize}
    \item means preserving partial operational capability under failures.
\end{itemize}

Examples:
\begin{itemize}
    \item reduced observation coverage;
    \item delayed communication;
    \item degraded coordination;
    \item simplified planning.
\end{itemize}

\begin{keybox}
    A resilient system does not necessarily preserve full performance. It preserves useful mission capability.
\end{keybox}

% ---------------------------------------------------------
% DISTRIBUTED RESILIENCE
% ---------------------------------------------------------

\subsection{Distributed Resilience}

Centralized systems are vulnerable because:
\begin{itemize}
    \item single failures may affect the entire system;
    \item communication disruptions may isolate agents.
\end{itemize}

Distributed autonomy improves resilience because:
\begin{itemize}
    \item local agents continue operating;
    \item coordination may adapt dynamically;
    \item failures remain partially localized.
\end{itemize}

Examples:
\begin{itemize}
    \item decentralized relaying;
    \item adaptive task reallocation;
    \item local fallback autonomy.
\end{itemize}

% ---------------------------------------------------------
% SELF-HEALING
% ---------------------------------------------------------

\subsection{Self-Healing Systems}

Self-healing systems:
\begin{itemize}
    \item detect failures autonomously;
    \item reorganize dynamically;
    \item redistribute responsibilities.
\end{itemize}

Examples:
\begin{itemize}
    \item replacing failed relay satellites;
    \item rerouting communication flows;
    \item redistributing observation tasks.
\end{itemize}

This strongly connects with:
\begin{itemize}
    \item organizational MAS;
    \item adaptive coordination;
    \item resilient scheduling.
\end{itemize}

% ---------------------------------------------------------
% RESILIENT MARL
% ---------------------------------------------------------

\subsection{Resilient MARL}

Interesting research directions:
\begin{itemize}
    \item robust MARL;
    \item adversarial MARL;
    \item fault-aware policies;
    \item uncertainty-aware coordination;
    \item adaptive communication strategies.
\end{itemize}

Key challenge:
\begin{itemize}
    \item ensuring stability under unexpected conditions.
\end{itemize}

% ---------------------------------------------------------
% OPERATIONAL PERSPECTIVE
% ---------------------------------------------------------

\subsection{Operational Perspective}

Operational operators care about:
\begin{itemize}
    \item mission continuity;
    \item predictability;
    \item safety;
    \item controllability.
\end{itemize}

This explains why:
\begin{itemize}
    \item explainability;
    \item organizational constraints;
    \item supervision;
\end{itemize}

are becoming increasingly important.

% ---------------------------------------------------------
% TYPICAL DISCUSSION
% ---------------------------------------------------------

\subsection{Typical Scientific Discussion}

\begin{questionbox}
    \textbf{Q: What is the difference between robustness and resilience?}

    A: Robustness means resisting disturbances without changing behavior. Resilience means adapting and recovering while preserving mission capability.
\end{questionbox}

\begin{questionbox}
    \textbf{Q: Why is distributed autonomy useful for resilience?}

    A: Because local agents can continue operating even when communication or centralized coordination is degraded.
\end{questionbox}

% \newpage

% =========================================================
% LESSON 2
% =========================================================

\lesson
{Neuro-Symbolic and Organizational Space Autonomy}
{Understand how symbolic structures can improve controllability and explainability in autonomous systems.}
{Pure learning-based systems often lack operational guarantees and interpretable coordination semantics.}

% ---------------------------------------------------------
% WHY NEURO-SYMBOLIC
% ---------------------------------------------------------

\subsection{Why Neuro-Symbolic Approaches Matter}

Pure machine learning systems:
\begin{itemize}
    \item can adapt well;
    \item but may remain difficult to interpret and control.
\end{itemize}

Pure symbolic systems:
\begin{itemize}
    \item are explainable;
    \item but often brittle and difficult to scale.
\end{itemize}

Neuro-symbolic approaches attempt to combine:
\begin{itemize}
    \item adaptation;
    \item structure;
    \item controllability;
    \item explainability.
\end{itemize}

\begin{keybox}
    The central idea is to combine learning-based adaptation with explicit operational knowledge.
\end{keybox}

% ---------------------------------------------------------
% ORGANIZATIONAL MODELS
% ---------------------------------------------------------

\subsection{Organizational Models}

Organizational approaches introduce:
\begin{itemize}
    \item roles;
    \item missions;
    \item norms;
    \item coordination constraints.
\end{itemize}

Examples in satellite constellations:
\begin{itemize}
    \item observer satellites;
    \item relay satellites;
    \item safety-oriented satellites;
    \item communication coordinators.
\end{itemize}

Benefits:
\begin{itemize}
    \item explicit operational doctrine;
    \item interpretable coordination;
    \item structured adaptation.
\end{itemize}

% ---------------------------------------------------------
% MOISE+
% ---------------------------------------------------------

\subsection{MOISE+ Organizational Modeling}

MOISE+ introduces three layers.

% ---------------------------------------------------------

\subsubsection*{1. Structural Layer}

Defines:
\begin{itemize}
    \item roles;
    \item relationships;
    \item organizational structure.
\end{itemize}

% ---------------------------------------------------------

\subsubsection*{2. Functional Layer}

Defines:
\begin{itemize}
    \item goals;
    \item missions;
    \item task decomposition.
\end{itemize}

% ---------------------------------------------------------

\subsubsection*{3. Deontic Layer}

Defines:
\begin{itemize}
    \item permissions;
    \item obligations;
    \item constraints.
\end{itemize}

% ---------------------------------------------------------
% ORGANIZATIONAL MARL
% ---------------------------------------------------------

\subsection{Organizational MARL}

Organizational MARL injects organizational semantics into learning.

Examples:
\begin{itemize}
    \item role-action constraints;
    \item role-reward guides;
    \item mission-level shaping;
    \item policy specialization.
\end{itemize}

Advantages:
\begin{itemize}
    \item improved controllability;
    \item more stable coordination;
    \item interpretable behaviors;
    \item operational alignment.
\end{itemize}

% ---------------------------------------------------------
% WORLD MODELS
% ---------------------------------------------------------

\subsection{World Models}

World models:
\begin{itemize}
    \item learn predictive models of the environment.
\end{itemize}

Applications:
\begin{itemize}
    \item planning;
    \item simulation;
    \item anomaly prediction;
    \item uncertainty estimation.
\end{itemize}

Interesting research direction:
\begin{itemize}
    \item symbolic constraints over learned world models.
\end{itemize}

% ---------------------------------------------------------
% EXPLAINABILITY
% ---------------------------------------------------------

\subsection{Trajectory-Level Explainability}

Modern research increasingly studies:
\begin{itemize}
    \item trajectory analysis;
    \item role inference;
    \item organizational pattern discovery;
    \item policy auditability.
\end{itemize}

Examples:
\begin{itemize}
    \item clustering behaviors;
    \item identifying emergent roles;
    \item extracting mission strategies.
\end{itemize}

% ---------------------------------------------------------
% CONNECTION WITH SPACE
% ---------------------------------------------------------

\subsection{Why this is Relevant for Space Systems}

Space operations require:
\begin{itemize}
    \item trust;
    \item predictability;
    \item safety;
    \item certification potential;
    \item operator understanding.
\end{itemize}

This explains why:
\begin{itemize}
    \item explainability;
    \item organizational structures;
    \item supervised autonomy;
\end{itemize}

are increasingly important.

% ---------------------------------------------------------
% TYPICAL DISCUSSION
% ---------------------------------------------------------

\subsection{Typical Scientific Discussion}

\begin{questionbox}
    \textbf{Q: Why use organizational constraints instead of reward shaping alone?}

    A: Because reward shaping does not necessarily enforce interpretable coordination semantics or operational discipline.
\end{questionbox}

\begin{questionbox}
    \textbf{Q: Why are neuro-symbolic approaches interesting for space autonomy?}

    A: Because they combine adaptability with explicit operational structure and improved explainability.
\end{questionbox}

% \newpage

% =========================================================
% LESSON 3
% =========================================================

\lesson
{Human-Autonomy Teaming}
{Understand how humans and autonomous systems can cooperate operationally.}
{Future missions are likely to rely on supervised autonomy rather than full autonomous replacement.}

% ---------------------------------------------------------
% WHY HUMAN-AUTONOMY TEAMING
% ---------------------------------------------------------

\subsection{Why Human-Autonomy Teaming Matters}

Space missions are:
\begin{itemize}
    \item safety-critical;
    \item mission-critical;
    \item politically sensitive;
    \item operationally expensive.
\end{itemize}

Operators therefore require:
\begin{itemize}
    \item visibility;
    \item control;
    \item explainability;
    \item trust.
\end{itemize}

\begin{keybox}
    The operational objective is often not full autonomy, but effective supervised autonomy.
\end{keybox}

% ---------------------------------------------------------
% OPERATOR ROLES
% ---------------------------------------------------------

\subsection{Role of Human Operators}

Operators typically:
\begin{itemize}
    \item define priorities;
    \item supervise execution;
    \item validate critical actions;
    \item manage anomalies;
    \item monitor mission health.
\end{itemize}

Autonomous systems:
\begin{itemize}
    \item support operators;
    \item reduce workload;
    \item react faster locally;
    \item manage routine coordination.
\end{itemize}

% ---------------------------------------------------------
% MIXED INITIATIVE
% ---------------------------------------------------------

\subsection{Mixed-Initiative Systems}

Mixed-initiative systems:
\begin{itemize}
    \item allow both humans and AI systems to influence planning and decisions.
\end{itemize}

Examples:
\begin{itemize}
    \item AI proposes plans;
    \item operators validate;
    \item operators override behaviors;
    \item autonomy handles execution details.
\end{itemize}

% ---------------------------------------------------------
% ADJUSTABLE AUTONOMY
% ---------------------------------------------------------

\subsection{Adjustable Autonomy}

Autonomy level may vary dynamically depending on:
\begin{itemize}
    \item operational context;
    \item communication availability;
    \item mission criticality;
    \item operator workload;
    \item environmental uncertainty.
\end{itemize}

Example:
\begin{itemize}
    \item low autonomy during nominal operations;
    \item high autonomy during communication blackout.
\end{itemize}

% ---------------------------------------------------------
% TRUST
% ---------------------------------------------------------

\subsection{Operator Trust}

Operators trust systems more when:
\begin{itemize}
    \item decisions are explainable;
    \item behaviors remain predictable;
    \item policies are auditable;
    \item operational doctrine is respected.
\end{itemize}

This strongly motivates:
\begin{itemize}
    \item explainable autonomy;
    \item organizational constraints;
    \item mission-level transparency.
\end{itemize}

% ---------------------------------------------------------
% EXPLAINABILITY
% ---------------------------------------------------------

\subsection{Explainability in Human-Autonomy Teaming}

Operators may ask:
\begin{itemize}
    \item Why was this decision taken?
    \item Why did coordination fail?
    \item Which role did the satellite assume?
    \item Why was this task prioritized?
\end{itemize}

Explainability therefore becomes:
\begin{itemize}
    \item operationally important;
    \item not only academically interesting.
\end{itemize}

% ---------------------------------------------------------
% CONNECTION WITH YOUR RESEARCH
% ---------------------------------------------------------

\subsection{Connection with Organizational MARL}

Organizational approaches naturally support:
\begin{itemize}
    \item interpretable coordination;
    \item role-based explanations;
    \item mission-level semantics;
    \item operator understanding.
\end{itemize}

Trajectory-level explainability methods:
\begin{itemize}
    \item can help audit learned policies;
    \item identify emergent organizational patterns;
    \item improve operator trust.
\end{itemize}

% ---------------------------------------------------------
% TYPICAL DISCUSSION
% ---------------------------------------------------------

\subsection{Typical Scientific Discussion}

\begin{questionbox}
    \textbf{Q: Why is human supervision still important in future autonomous missions?}

    A: Because operators remain responsible for mission safety, strategic objectives, and operational accountability.
\end{questionbox}

\begin{questionbox}
    \textbf{Q: Why is explainability operationally important?}

    A: Because operators must understand and trust autonomous decisions before deploying them in mission-critical systems.
\end{questionbox}

% \newpage

% =========================================================
% LESSON 4
% =========================================================

\lesson
{Frontier Topics and Emerging Directions}
{Understand the main emerging research topics likely to shape future MASSpace discussions.}
{These directions represent current and future challenges in scalable space autonomy.}

% ---------------------------------------------------------
% MEGA CONSTELLATIONS
% ---------------------------------------------------------

\subsection{Mega-Constellations}

Mega-constellations involve:
\begin{itemize}
    \item hundreds or thousands of satellites.
\end{itemize}

Challenges:
\begin{itemize}
    \item scalability;
    \item distributed scheduling;
    \item communication management;
    \item collision risk;
    \item autonomous coordination.
\end{itemize}

Examples:
\begin{itemize}
    \item Starlink;
    \item OneWeb;
    \item Kuiper.
\end{itemize}

\begin{keybox}
    Mega-constellations strongly motivate distributed AI and scalable autonomy.
\end{keybox}

% ---------------------------------------------------------
% SPACE TRAFFIC MANAGEMENT
% ---------------------------------------------------------

\subsection{Space Traffic Management (STM)}

STM focuses on:
\begin{itemize}
    \item orbital congestion;
    \item collision avoidance;
    \item conjunction management;
    \item coordination between operators.
\end{itemize}

Important because:
\begin{itemize}
    \item LEO congestion is increasing rapidly.
\end{itemize}

Interesting research topics:
\begin{itemize}
    \item decentralized collision avoidance;
    \item cooperative maneuver planning;
    \item autonomous conjunction assessment.
\end{itemize}

% ---------------------------------------------------------
% DIGITAL TWINS
% ---------------------------------------------------------

\subsection{Digital Twins}

Digital twins:
\begin{itemize}
    \item are high-fidelity virtual replicas of physical systems.
\end{itemize}

Applications:
\begin{itemize}
    \item simulation;
    \item predictive maintenance;
    \item mission rehearsal;
    \item anomaly diagnosis;
    \item autonomy validation.
\end{itemize}

Potential connection:
\begin{itemize}
    \item world models for operational autonomy.
\end{itemize}

% ---------------------------------------------------------
% AI-DRIVEN MISSION CONTROL
% ---------------------------------------------------------

\subsection{AI-Driven Mission Control}

Future mission control systems may increasingly integrate:
\begin{itemize}
    \item AI-assisted planning;
    \item anomaly detection;
    \item autonomous scheduling;
    \item operator-support systems.
\end{itemize}

Goal:
\begin{itemize}
    \item reduce operator workload;
    \item improve scalability;
    \item accelerate response time.
\end{itemize}

% ---------------------------------------------------------
% DEEP SPACE AUTONOMY
% ---------------------------------------------------------

\subsection{Deep-Space Autonomy}

Deep-space missions involve:
\begin{itemize}
    \item very large communication delays;
    \item isolation;
    \item uncertain environments.
\end{itemize}

Examples:
\begin{itemize}
    \item Mars exploration;
    \item lunar systems;
    \item asteroid missions.
\end{itemize}

This strongly motivates:
\begin{itemize}
    \item onboard planning;
    \item local autonomy;
    \item adaptive mission management.
\end{itemize}

% ---------------------------------------------------------
% SPACE SWARMS
% ---------------------------------------------------------

\subsection{Space Swarms}

Future systems may involve:
\begin{itemize}
    \item large autonomous satellite swarms;
    \item collaborative sensing;
    \item decentralized mission orchestration.
\end{itemize}

Potential applications:
\begin{itemize}
    \item distributed Earth observation;
    \item adaptive communication networks;
    \item resilient sensing systems.
\end{itemize}

Interesting research directions:
\begin{itemize}
    \item emergent coordination;
    \item organizational adaptation;
    \item swarm-level explainability.
\end{itemize}

% ---------------------------------------------------------
% SPACE AI CHALLENGES
% ---------------------------------------------------------

\subsection{Open Challenges in Space AI}

Major open problems:
\begin{itemize}
    \item safe autonomy;
    \item explainability;
    \item scalability;
    \item robustness;
    \item certification;
    \item human trust;
    \item operational deployment.
\end{itemize}

\begin{warningbox}
    The main challenge is no longer only achieving autonomy, but achieving trustworthy operational autonomy.
\end{warningbox}

% ---------------------------------------------------------
% RESEARCH POSITIONING
% ---------------------------------------------------------

\subsection{Research Positioning}

Your research profile is particularly aligned with:
\begin{itemize}
    \item organizational autonomy;
    \item resilient MARL;
    \item neuro-symbolic coordination;
    \item explainable autonomy;
    \item cyber-resilient space systems.
\end{itemize}

These intersections are:
\begin{itemize}
    \item relatively uncommon;
    \item highly relevant;
    \item increasingly important operationally.
\end{itemize}

% ---------------------------------------------------------
% TYPICAL DISCUSSION
% ---------------------------------------------------------

\subsection{Typical Scientific Discussion}

\begin{questionbox}
    \textbf{Q: Why are organizational approaches promising for future space autonomy?}

    A: Because they can improve controllability, explainability, and operational alignment in large distributed autonomous systems.
\end{questionbox}

\begin{questionbox}
    \textbf{Q: What is one of the main future challenges for space AI?}

    A: Achieving scalable and trustworthy autonomy under degraded and uncertain operational conditions.
\end{questionbox}


% =========================================================
% APPENDICES
% =========================================================

\appendix

\part{Appendices}

% =========================================================
% APPENDIX 1
% =========================================================

\section{Core English Glossary}

This glossary contains the most important technical vocabulary likely to appear:
\begin{itemize}
    \item during MASSpace presentations;
    \item in space autonomy papers;
    \item in workshop discussions;
    \item during networking conversations.
\end{itemize}

The objective is not perfect memorization, but:
\begin{itemize}
    \item fast recognition;
    \item operational understanding;
    \item fluency during discussions.
\end{itemize}

% ---------------------------------------------------------
% GENERAL SPACE
% ---------------------------------------------------------

\subsection{General Space Vocabulary}

\begin{tabularx}{\columnwidth}{lX}
    \toprule
    Term               & Meaning                                                   \\
    \midrule
    Space segment      & Satellites and spacecraft in orbit                        \\
    Ground segment     & Ground infrastructure and mission control                 \\
    Payload            & Mission-specific instrument or subsystem                  \\
    Constellation      & Coordinated group of satellites                           \\
    Mission operations & Operational management of spacecraft and missions         \\
    Spacecraft         & General term for a vehicle operating in space             \\
    Orbital pass       & Satellite flying over a region/station                    \\
    Ground track       & Orbit projection on Earth surface                         \\
    Line of sight      & Direct visibility between entities                        \\
    Visibility window  & Time interval where communication/observation is possible \\
    \bottomrule
\end{tabularx}

% ---------------------------------------------------------
% ORBITS
% ---------------------------------------------------------

\subsection{Orbital Mechanics Vocabulary}

\begin{tabularx}{\columnwidth}{lX}
    \toprule
    Term                         & Meaning                                                  \\
    \midrule
    LEO                          & Low Earth Orbit                                          \\
    MEO                          & Medium Earth Orbit                                       \\
    GEO                          & Geostationary Orbit                                      \\
    Inclination                  & Orbit tilt relative to equator                           \\
    Orbital period               & Time required to complete one orbit                      \\
    Sun-synchronous orbit        & Orbit maintaining constant solar illumination conditions \\
    Conjunction                  & Close approach between orbital objects                   \\
    Collision avoidance maneuver & Orbit adjustment reducing collision risk                 \\
    SSA                          & Space Situational Awareness                              \\
    SDA                          & Space Domain Awareness                                   \\
    \bottomrule
\end{tabularx}

% ---------------------------------------------------------
% COMMUNICATIONS
% ---------------------------------------------------------

\subsection{Communication Vocabulary}

\begin{tabularx}{\columnwidth}{lX}
    \toprule
    Term                      & Meaning                                              \\
    \midrule
    Uplink                    & Ground-to-satellite communication                    \\
    Downlink                  & Satellite-to-ground communication                    \\
    ISL / Crosslink           & Inter-satellite communication                        \\
    Bandwidth                 & Data transmission capacity                           \\
    Latency                   & Communication delay                                  \\
    Communication window      & Time interval during which communication is possible \\
    DTN                       & Delay-Tolerant Networking                            \\
    Intermittent connectivity & Non-continuous communication availability            \\
    Relay                     & Forwarding data through intermediate satellite       \\
    \bottomrule
\end{tabularx}

% ---------------------------------------------------------
% AUTONOMY
% ---------------------------------------------------------

\subsection{Autonomy and AI Vocabulary}

\begin{tabularx}{\columnwidth}{lX}
    \toprule
    Term                            & Meaning                                              \\
    \midrule
    Onboard autonomy                & Autonomous decision-making executed onboard          \\
    Distributed autonomy            & Decentralized decision-making across agents          \\
    Replanning                      & Dynamic modification of plans                        \\
    Mission continuity              & Maintaining mission capability under disruptions     \\
    Graceful degradation            & Preserving partial capability under failures         \\
    Operational assist              & AI supporting operators operationally                \\
    Human-in-the-loop               & Human supervision retained in autonomy loop          \\
    Explainable autonomy            & Interpretable autonomous decision-making             \\
    Trajectory-level explainability & Explaining behaviors using trajectories and policies \\
    \bottomrule
\end{tabularx}

% ---------------------------------------------------------
% MAS / MARL
% ---------------------------------------------------------

\subsection{MAS and MARL Vocabulary}

\begin{tabularx}{\columnwidth}{lX}
    \toprule
    Term                       & Meaning                                                    \\
    \midrule
    Dec-POMDP                  & Decentralized Partially Observable Markov Decision Process \\
    CTDE                       & Centralized Training, Decentralized Execution              \\
    Task allocation            & Assigning tasks to agents                                  \\
    Distributed coordination   & Coordinated decentralized behavior                         \\
    Emergent behavior          & Coordination not explicitly programmed                     \\
    Role specialization        & Agents adopting distinct functions                         \\
    Organizational constraints & Explicit coordination structure                            \\
    Role-action guide          & Organizational restriction on actions                      \\
    Mission-aware policy       & Policy aligned with mission semantics                      \\
    \bottomrule
\end{tabularx}

% ---------------------------------------------------------
% CYBERSECURITY
% ---------------------------------------------------------

\subsection{Cybersecurity and Resilience Vocabulary}

\begin{tabularx}{\columnwidth}{lX}
    \toprule
    Term                  & Meaning                                     \\
    \midrule
    Jamming               & Communication interference                  \\
    Spoofing              & Fake signal/data injection                  \\
    Contested environment & Adversarial/degraded environment            \\
    Fault tolerance       & Ability to continue despite faults          \\
    Cyber-resilience      & Ability to resist and recover from attacks  \\
    Adaptive autonomy     & Dynamically adapting autonomous behavior    \\
    Mission assurance     & Ensuring mission reliability and continuity \\
    Self-healing system   & System reorganizing after failures          \\
    \bottomrule
\end{tabularx}

% \newpage

% =========================================================
% APPENDIX 2
% =========================================================

\section{Typical MASSpace Expressions}

This section contains expressions commonly used in:
\begin{itemize}
    \item presentations;
    \item papers;
    \item discussions;
    \item Q\&A sessions.
\end{itemize}

Understanding these formulations is extremely important for:
\begin{itemize}
    \item fluency;
    \item networking;
    \item scientific discussions.
\end{itemize}

% ---------------------------------------------------------
% COMMON EXPRESSIONS
% ---------------------------------------------------------

\subsection{Very Common Expressions}

% ---------------------------------------------------------

\subsubsection*{``Mission-critical''}

Meaning:
\begin{itemize}
    \item essential for mission success.
\end{itemize}

Example:
\begin{quote}
    ``This is a mission-critical operation.''
\end{quote}

% ---------------------------------------------------------

\subsubsection*{``Degraded conditions''}

Meaning:
\begin{itemize}
    \item partially disrupted operational context.
\end{itemize}

Example:
\begin{quote}
    ``The system must remain operational under degraded conditions.''
\end{quote}

% ---------------------------------------------------------

\subsubsection*{``Contested environment''}

Meaning:
\begin{itemize}
    \item adversarial or hostile environment.
\end{itemize}

Example:
\begin{quote}
    ``Future space systems may operate in contested environments.''
\end{quote}

% ---------------------------------------------------------

\subsubsection*{``Scalable autonomy''}

Meaning:
\begin{itemize}
    \item autonomy capable of handling large-scale systems.
\end{itemize}

Example:
\begin{quote}
    ``Mega-constellations require scalable autonomy.''
\end{quote}

% ---------------------------------------------------------

\subsubsection*{``Operationally relevant''}

Meaning:
\begin{itemize}
    \item realistic and useful for operations.
\end{itemize}

Example:
\begin{quote}
    ``We focus on operationally relevant scenarios.''
\end{quote}

% ---------------------------------------------------------

\subsubsection*{``Human-in-the-loop''}

Meaning:
\begin{itemize}
    \item humans remain involved in decision-making.
\end{itemize}

Example:
\begin{quote}
    ``The operator remains in the loop.''
\end{quote}

% ---------------------------------------------------------

\subsubsection*{``Distributed decision-making''}

Meaning:
\begin{itemize}
    \item decentralized autonomous decisions.
\end{itemize}

Example:
\begin{quote}
    ``Distributed decision-making is required under communication disruptions.''
\end{quote}

% ---------------------------------------------------------

\subsubsection*{``Mission continuity''}

Meaning:
\begin{itemize}
    \item preserving mission capability despite disruptions.
\end{itemize}

Example:
\begin{quote}
    ``The system prioritizes mission continuity under failures.''
\end{quote}

% ---------------------------------------------------------

\subsubsection*{``Graceful degradation''}

Meaning:
\begin{itemize}
    \item reduced but still functional operations.
\end{itemize}

Example:
\begin{quote}
    ``The architecture supports graceful degradation.''
\end{quote}

% ---------------------------------------------------------

\subsubsection*{``Supervised autonomy''}

Meaning:
\begin{itemize}
    \item autonomous execution under human supervision.
\end{itemize}

Example:
\begin{quote}
    ``The goal is supervised autonomy rather than full autonomous replacement.''
\end{quote}

% \newpage

% =========================================================
% APPENDIX 3
% =========================================================

\section{Typical MASSpace Questions and Answers}

This section contains:
\begin{itemize}
    \item common workshop questions;
    \item expected scientific reasoning;
    \item safe operational framings.
\end{itemize}

The goal is not memorization:
\begin{itemize}
    \item but discussion readiness.
\end{itemize}

% ---------------------------------------------------------
% REALISM
% ---------------------------------------------------------

\subsection{Simulation Realism}

\begin{questionbox}
    \textbf{Q: How realistic is your simulator?}

    A: The objective is not flight-grade orbital realism, but controlled experimentation on operational coordination, communication constraints, resilience, and distributed autonomy.
\end{questionbox}

% ---------------------------------------------------------
% WHY MARL
% ---------------------------------------------------------

\subsection{Why MARL?}

\begin{questionbox}
    \textbf{Q: Why use MARL for satellite constellations?}

    A: Because constellations naturally involve decentralized agents operating under partial observability, communication constraints, and shared mission objectives.
\end{questionbox}

% ---------------------------------------------------------
% WHY NOT PLANNING
% ---------------------------------------------------------

\subsection{Why Not Pure Planning?}

\begin{questionbox}
    \textbf{Q: Why not rely only on centralized planning or DCOP approaches?}

    A: Centralized planning remains extremely valuable, but future constellations may require more adaptive decentralized autonomy under communication disruptions and scalability constraints.
\end{questionbox}

% ---------------------------------------------------------
% WHY ORGANIZATIONAL
% ---------------------------------------------------------

\subsection{Why Organizational Constraints?}

\begin{questionbox}
    \textbf{Q: Why use organizational constraints instead of reward shaping alone?}

    A: Because reward shaping does not necessarily enforce interpretable operational semantics, explicit role specialization, or mission-aware coordination structures.
\end{questionbox}

% ---------------------------------------------------------
% HUMAN SUPERVISION
% ---------------------------------------------------------

\subsection{Human Supervision}

\begin{questionbox}
    \textbf{Q: Do you envision fully autonomous missions?}

    A: The operational objective is usually supervised autonomy where operators remain responsible for strategic decisions and mission validation.
\end{questionbox}

% ---------------------------------------------------------
% EXPLAINABILITY
% ---------------------------------------------------------

\subsection{Explainability}

\begin{questionbox}
    \textbf{Q: Why is explainability important in space systems?}

    A: Because operators require trust, auditability, predictability, and operational understanding before deploying autonomous systems in mission-critical environments.
\end{questionbox}

% ---------------------------------------------------------
% RESILIENCE
% ---------------------------------------------------------

\subsection{Resilience}

\begin{questionbox}
    \textbf{Q: Why is resilience central in future constellations?}

    A: Because large distributed systems inevitably experience failures, communication disruptions, environmental hazards, and potentially adversarial conditions.
\end{questionbox}

% \newpage

% =========================================================
% APPENDIX 4
% =========================================================

\section{Major Organizations and Ecosystem}

Knowing the major actors helps:
\begin{itemize}
    \item contextualize discussions;
    \item understand references;
    \item identify collaboration ecosystems.
\end{itemize}

% ---------------------------------------------------------
% SPACE AGENCIES
% ---------------------------------------------------------

\subsection{Major Space Agencies}

\begin{tabularx}{\columnwidth}{lX}
    \toprule
    Organization & Role                                                     \\
    \midrule
    NASA         & US space agency; exploration, science, autonomy research \\
    ESA          & European Space Agency                                    \\
    CNES         & French space agency                                      \\
    DLR          & German aerospace research center                         \\
    JAXA         & Japanese space agency                                    \\
    ISRO         & Indian space agency                                      \\
    CSA          & Canadian Space Agency                                    \\
    \bottomrule
\end{tabularx}

% ---------------------------------------------------------
% RESEARCH ORGANIZATIONS
% ---------------------------------------------------------

\subsection{Research and Defense Organizations}

\begin{tabularx}{\columnwidth}{lX}
    \toprule
    Organization               & Role                                    \\
    \midrule
    ONERA                      & French aerospace research organization  \\
    DARPA                      & US advanced defense research agency     \\
    AFRL                       & US Air Force Research Laboratory        \\
    ESA Advanced Concepts Team & Emerging space technologies and AI      \\
    NASA JPL                   & Robotics, autonomy, exploration systems \\
    \bottomrule
\end{tabularx}

% ---------------------------------------------------------
% INDUSTRIAL ACTORS
% ---------------------------------------------------------

\subsection{Industrial Ecosystem}

\begin{tabularx}{\columnwidth}{lX}
    \toprule
    Company                  & Main focus                             \\
    \midrule
    SpaceX                   & Launch systems and mega-constellations \\
    Airbus Defence and Space & European aerospace systems             \\
    Thales Alenia Space      & Satellites and space infrastructures   \\
    Planet Labs              & Earth observation constellations       \\
    ICEYE                    & Radar imaging satellites               \\
    Maxar                    & High-resolution Earth observation      \\
    OneWeb                   & Satellite communications constellation \\
    \bottomrule
\end{tabularx}

% \newpage

% =========================================================
% APPENDIX 5
% =========================================================

\section{Research Positioning Notes}

This section summarizes how your profile aligns with MASSpace topics.

% ---------------------------------------------------------
% YOUR STRENGTHS
% ---------------------------------------------------------

\subsection{Your Main Strengths}

Your profile is particularly aligned with:
\begin{itemize}
    \item organizational MARL;
    \item resilient distributed autonomy;
    \item explainable coordination;
    \item neuro-symbolic systems;
    \item cyber-resilient autonomy.
\end{itemize}

These intersections are:
\begin{itemize}
    \item relatively uncommon;
    \item scientifically interesting;
    \item operationally relevant.
\end{itemize}

% ---------------------------------------------------------
% IMPORTANT POSITIONING
% ---------------------------------------------------------

\subsection{Important Positioning Strategy}

Do NOT try to position yourself as:
\begin{center}
    ``orbital mechanics expert''
\end{center}

Instead, position yourself as:
\begin{center}
    ``AI/autonomy/MAS researcher addressing future operational challenges in distributed space systems''
\end{center}

This is:
\begin{itemize}
    \item accurate;
    \item credible;
    \item strategically strong.
\end{itemize}

% ---------------------------------------------------------
% IMPORTANT FRAMINGS
% ---------------------------------------------------------

\subsection{Strong Scientific Framings}

Good framings:
\begin{itemize}
    \item resilient distributed autonomy;
    \item mission-aware coordination;
    \item supervised autonomy;
    \item organizational guidance;
    \item operational explainability;
    \item scalable constellation management.
\end{itemize}

Avoid:
\begin{itemize}
    \item overclaiming realism;
    \item claiming fully autonomous replacement;
    \item overselling deployment readiness.
\end{itemize}

% ---------------------------------------------------------
% NETWORKING
% ---------------------------------------------------------

\subsection{Networking Advice}

During discussions:
\begin{itemize}
    \item ask operational questions;
    \item discuss scalability challenges;
    \item ask about communication disruptions;
    \item discuss explainability and supervision;
    \item connect MAS concepts with operational problems.
\end{itemize}

Good discussion topics:
\begin{itemize}
    \item decentralized scheduling;
    \item resilient coordination;
    \item explainable autonomy;
    \item organizational structures;
    \item human-autonomy teaming.
\end{itemize}

% ---------------------------------------------------------
% FINAL NOTE
% ---------------------------------------------------------

\begin{keybox}
    You do not need to become a pure aerospace engineer.
    Your strongest positioning is becoming an expert in resilient distributed AI/autonomy for future space systems.
\end{keybox}

\end{document}